\documentclass[11pt,a4paper]{article}
\usepackage[margin=26mm]{geometry}
\usepackage[T1]{fontenc}
\usepackage{lmodern,microtype,amsmath,amssymb,amsthm,mathtools}
\usepackage{enumitem,needspace,xcolor,hyperref}
\definecolor{ink}{HTML}{16364A}
\hypersetup{colorlinks=true,linkcolor=ink,urlcolor=ink,citecolor=ink,
pdfauthor={Zhuoni Chi and Shaozhen Cao},
pdftitle={Distribution of Singular Data Generated by Compact Forced Navier--Stokes Blowup}}
\newtheorem{theorem}{Theorem}[section]
\newtheorem{proposition}[theorem]{Proposition}
\newtheorem{lemma}[theorem]{Lemma}
\newtheorem{corollary}[theorem]{Corollary}
\theoremstyle{definition}\newtheorem{definition}[theorem]{Definition}
\theoremstyle{remark}\newtheorem{remark}[theorem]{Remark}
\newcommand{\R}{\mathbb R}\newcommand{\TT}{\mathbb T^3}
\newcommand{\FF}{\mathcal F}\newcommand{\XX}{\mathcal X}
\newcommand{\BB}{\mathcal B}\newcommand{\PP}{\mathbb P}
\newcommand{\dd}{\,\mathrm d}\newcommand{\eps}{\varepsilon}
\newcommand{\norm}[1]{\left\lVert#1\right\rVert}
\newcommand{\ip}[2]{\left\langle#1,#2\right\rangle}
\DeclareMathOperator{\supp}{supp}
\title{Distribution of Singular Data Generated by\\ Compact Forced Navier--Stokes Blowup}
\author{Zhuoni Chi\\\texttt{znchi@zju.edu.cn}\and Shaozhen Cao\\\texttt{shaozhencao@zju.edu.cn}}
\date{9 September 2026}
\begin{document}
\maketitle
\begin{abstract}
Starting from the compact, smoothly forced Navier--Stokes blowup solution constructed by OpenAI, we study the distribution of the singular data it generates. For fixed viscosity and time horizon on the three-dimensional torus, smooth forces producing classical breakdown from rest by that time are dense in the inherited time-integrated spatial $H^s$ topology if and only if $s<1/2$ (the norms are defined in Section~\ref{sec:intro}). The positive result follows from an exact insertion into every regular reference trajectory. A localized vector potential removes the background around a concentrated singular packet, so all nonlinear cross terms vanish and the modified force remains smooth through the singular time. We give the full support construction, derivative estimates, fractional Sobolev scaling and classical-lifespan argument. The inserted trajectories converge strongly in the energy and dissipation norm. A separate critical-force bootstrap, followed by an $H^1$ estimate and high-Sobolev continuation, supplies the regular open set needed for non-density at and above $s=1/2$. Further results describe extended-data projections, every fixed smooth initial-velocity slice, mixed force norms, infinite-dimensional variations and interior no-slip realizations. The construction varies the force; it does not classify singular initial velocities for a single prescribed force.
\end{abstract}
\noindent\textbf{Keywords:} forced Navier--Stokes equations; singular data; localized insertion; critical Sobolev topology; energy convergence.

\section{Introduction and main result}\label{sec:intro}
The initial datum of a forced evolution equation is only one part of its input. The external force is another. Accordingly, the distribution of singular inputs can mean the distribution of initial velocities at fixed force, the distribution of forces at fixed initial velocity, or the distribution of pairs. A singular solution starting from rest has just one initial velocity, even if its force admits a large family of variations.

OpenAI's \emph{Finite Time Blowup for Navier--Stokes} \cite[Theorem 1.1]{openai} supplies, for every positive viscosity, a compact whole-space solution with zero initial velocity, globally smooth compactly supported forcing, bounded kinetic energy and unbounded speed at time one. We take this theorem and its associated Lean formalization as established results. The compact packet is restated in Theorem~\ref{thm:packet}. The present article proves distribution properties of the data obtained from that packet.

Hofmanov\'a, Zhu and Zhu \cite[Theorem~1.4, Corollary~4.6 and Theorem~5.1]{hzz2024} provide a close antecedent: forces producing nonunique Leray--Hopf solutions are dense in $L^1(0,1;L^2(\R^3))$ from rest, and every $L^2_\sigma$ initial velocity admits a force producing nonuniqueness. This motivates our emphasis on force topology and on fixing the initial velocity while varying the force. Here the target property is classical breakdown under globally smooth forcing, and the construction removes a regular background locally before inserting the compact packet. The weak-solution nonuniqueness theorem does not supply that packet or the sharp Sobolev thresholds proved here; our conclusions do not assert nonuniqueness of Leray--Hopf continuations.

The analytic prerequisites are elementary calculus, Lebesgue and Bochner integration, Fourier inversion and Plancherel, H\"older's inequality, and the contraction principle. The critical embeddings are proved in Appendix~\ref{sec:critical-appendix}; the multiplication estimates, local evolution and continuation arguments used here are proved in the text. The other references provide historical and conceptual context, rather than additional PDE existence inputs.

Fix $\nu>0$ and $T>0$. The principal domain is the unit torus
$\TT=\R^3/\mathbb Z^3$, with volume one.

For $s\in\R$ and a periodic scalar or vector distribution $z$, write
$\widehat z(k)=\int_{\TT}z(x)e^{-2\pi i k\cdot x}\dd x$, with Fourier coefficients understood distributionally when necessary. The inhomogeneous space $H^s(\TT)$ consists of distributions with finite norm
\[
 \norm{z}_{H^s}^2
 =\sum_{k\in\mathbb Z^3}(1+4\pi^2|k|^2)^s|\widehat z(k)|^2.
\]
The homogeneous space $\dot H^s(\TT)$ consists of mean-zero distributions with finite norm
\[
 \norm{z}_{\dot H^s}^2
 =\sum_{k\ne0}(2\pi|k|)^{2s}|\widehat z(k)|^2.
\]
The homogeneous and inhomogeneous norms are equivalent on mean-zero fields at each fixed order. Vector and tensor norms are the square root of the sum of the squared component norms. Unless a domain is shown, a spatial norm in the periodic argument is taken on $\TT$.

For a strongly measurable map $z:I\to X$ into a normed spatial space, use
\begin{equation}\label{eq:time-norms}
 \norm{z}_{L^q(I;X)}=
 \left(\int_I\norm{z(t)}_X^q\dd t\right)^{1/q}\quad(1\le q<\infty),
 \qquad
 \norm{z}_{L^\infty(I;X)}=\operatorname*{ess\,sup}_{t\in I}\norm{z(t)}_X.
\end{equation}
Thus $L^1_tH^s_x$ for forces means $L^1(0,\infty;H^s(\TT))$, with norm
$\int_0^\infty\norm{f(t)}_{H^s(\TT)}\dd t$; $H^s$ measures spatial regularity only.
Unless an interval is displayed, mixed force norms use $(0,\infty)$ and
presingular velocity norms use $(0,T)$. The same convention applies to
$X=L^p$ and $X=\dot H^s$, and to the explicitly named domains below.

We consider
\begin{equation}\label{eq:NS}
 \partial_tu+(u\cdot\nabla)u-\nu\Delta u+\nabla p=f,
 \qquad \nabla\cdot u=0,\qquad u(0)=a.
\end{equation}
Both $u$ and $p$ are periodic; pressure is normalized by $\int_{\TT}p=0$.
The smooth input spaces are
\begin{equation}\label{eq:inputspaces}
 \XX=C^\infty_{\mathrm{div}}(\TT;\R^3),\qquad
 \FF=C_c^\infty(\TT\times(0,\infty);\R^3).
\end{equation}
The support restriction in $\FF$ concerns positive time. In particular, each force is defined and smooth through any possible terminal time of its velocity.

Let $T_{\max}^{\nu}(a,f)$ denote the maximal classical lifespan for the pressure-normalized problem. Its existence and uniqueness framework is recalled in Proposition~\ref{prop:local}. Define
\begin{equation}\label{eq:singularforces}
 \BB_{\nu,a,T}=\{f\in\FF:T_{\max}^{\nu}(a,f)\le T\},
 \qquad \BB^0_{\nu,T}=\BB_{\nu,0,T}.
\end{equation}

\begin{theorem}[Sobolev density threshold]\label{thm:main}
Equip $\FF$ with its relative $L^1(0,\infty;H^s(\TT))$ norm topology.
\begin{enumerate}[label=(\roman*),nosep]
\item For every fixed $a\in\XX$, the set $\BB_{\nu,a,T}$ is dense if $s<1/2$.
\item On the zero-initial-velocity slice,
\[
 \BB^0_{\nu,T}\text{ is dense in }\FF
 \quad\Longleftrightarrow\quad s<\tfrac12.
\]
\end{enumerate}
\end{theorem}

Here density is relative to smooth forces in a specified norm. It is not density in the usual test-function topology, which imposes control of all derivatives on a common support. The approximating forces below concentrate in shrinking regions and have increasing peak amplitudes.

The construction is local in space and late in time. Given a reference solution regular through $T$, we preserve its initial velocity and its history until shortly before $T$, then insert a packet that becomes singular exactly at $T$. The resulting force is smooth on the entire positive time axis. The velocity perturbation tends to zero in
\begin{equation}\label{eq:Enorm}
 \norm{z}_{E_T}
 :=\norm{z}_{L^\infty(0,T;L^2)}
   +\norm{\nabla z}_{L^2(0,T;L^2)}.
\end{equation}
This is equivalent, on a fixed finite interval, to the sum norm of
$L^\infty_tL^2_x\cap L^2_tH^1_x$. It requires neither a classical value nor a continuation at $T$.

The proof separates three issues. First, parabolic concentration makes the packet small in subcritical force norms. Second, a divergence-free background correction makes the nonlinear insertion exact while keeping its force smooth at the singular endpoint. Third, a separate critical regularity estimate proves the existence of an open regular neighborhood. A scaling calculation alone would not establish non-density.

Section~\ref{sec:analytic} fixes the analytic conventions and the continuation criterion. Section~\ref{sec:packet} derives the packet estimates, including the fractional localization bound. Section~\ref{sec:insertion} proves the insertion theorem and its interior-boundary extension. Section~\ref{sec:density} treats the density and projection quantifiers. Section~\ref{sec:critical} establishes the critical obstruction and completes Theorem~\ref{thm:main}. Section~\ref{sec:variants} gives additional families.

\section{Analytic conventions and classical continuation}\label{sec:analytic}
\subsection{Sobolev norms, projection and pressure}
\paragraph{Euclidean and interior-domain norms.}
For a scalar distribution $Z$ on $\R^3$, our Fourier convention is
$\widehat Z(\xi)=(2\pi)^{-3/2}\int_{\R^3}e^{-ix\cdot\xi}Z(x)\dd x$,
initially on Schwartz functions and then by duality. For an open domain
$\Omega\subset\R^3$ and every $s\in\R$ set
\begin{equation}\label{eq:restriction-norm}
 \norm{Z}_{H^s(\R^3)}^2
 =\int_{\R^3}(1+|\xi|^2)^s|\widehat Z(\xi)|^2\dd\xi,
 \qquad
 \norm{z}_{H^s(\Omega)}
 =\inf_{Z|_\Omega=z}\norm{Z}_{H^s(\R^3)}.
\end{equation}
Here $H^s(\R^3)$ consists of tempered distributions with finite displayed
norm, and $H^s(\Omega)$ is its space of distributional restrictions with the
quotient norm. In particular, negative orders use this restriction definition,
not a separately chosen dual-space convention. At $s=0$ this is the usual
$L^2(\Omega)$ norm. The auxiliary Euclidean homogeneous norm is
$\norm{Z}_{\dot H^s(\R^3)}^2=\int_{\R^3}|\xi|^{2s}|\widehat Z(\xi)|^2\dd\xi$
whenever finite; below it is used only for nonnegative orders on compact smooth
profiles in the scaling argument. It does not replace the inhomogeneous norm
in the force topology.

For every fixed compact $K\Subset\Omega$ and $s\in\R$, smooth fields supported
in $K$ obey
\begin{equation}\label{eq:zero-extension}
 \norm{z}_{H^s(\Omega)}\le\norm{E_0z}_{H^s(\R^3)}
 \le C_{s,K,\Omega}\norm{z}_{H^s(\Omega)},
 \qquad \operatorname{supp}z\subset K,
\end{equation}
where $E_0z$ is extension by zero. To prove the second inequality, fix
$\chi\in C_c^\infty(\Omega)$ equal to one near $K$. For every extension $Z$
of $z$, $\chi Z=E_0z$. Multiplication by this fixed $\chi$ is bounded on
$H^s(\R^3)$ for every real $s$: the Fourier convolution formula and
\[
 \frac{(1+|\xi|^2)^{s/2}}{(1+|\eta|^2)^{s/2}}
 \le C_s(1+|\xi-\eta|^2)^{|s|/2}
\]
reduce the estimate to Young's $L^1*L^2\to L^2$ inequality, since
$\widehat\chi$ decays faster than every power. Taking the infimum over $Z$
proves the bound; the first inequality follows from the definition.
The constant depends on the fixed cutoff and $s$, not on a smaller support
or a shrinking scale $\eps$. Applying the pointwise bounds in time gives the
same comparison for all the mixed norms in \eqref{eq:time-norms}, provided
$K$ is fixed for all times. No bounded zero-extension operator on arbitrary
$H^s(\Omega)$ is asserted. All vector and tensor norms in these definitions
are the square root of the sum of the squared component norms.

The Leray projection $\PP$ is the Fourier multiplier
\[
 \widehat{\PP z}(k)=
 \left(I-\frac{k\otimes k}{|k|^2}\right)\widehat z(k)\quad(k\ne0),
 \qquad \widehat{\PP z}(0)=\widehat z(0).
\]
It is bounded on every $H^s$, commutes with spatial derivatives and the heat semigroup, and eliminates periodic gradients. Equation~\eqref{eq:NS} becomes
\begin{equation}\label{eq:projected}
 \partial_tu-\nu\Delta u
 =-\PP\nabla\cdot(u\otimes u)+\PP f.
\end{equation}
The pressure is then recovered uniquely by
\[
 \Delta p=\nabla\cdot f-\sum_{i,j}\partial_i\partial_j(u_i u_j),
 \qquad \int_{\TT}p=0.
\]
The right side has zero mean. Periodicity of the pressure is part of the solution class; a nonperiodic affine pressure is not an allowed alternative gauge.

\begin{lemma}[Sobolev estimates used below]\label{lem:calculus}
For each integer $k\ge3$ there are constants depending only on $k$ and the torus such that
\begin{align}
 \norm{zw}_{H^k}&\le C_k\norm{z}_{H^k}\norm{w}_{H^k},\label{eq:algebra}\\
 \norm{u\otimes u}_{H^k}
 &\le C_k\norm{u}_{H^2}\norm{u}_{H^k}.\label{eq:tame}
\end{align}
For mean-zero $v$,
\begin{equation}\label{eq:embeddings}
 \norm{v}_3\le C\norm{v}_{\dot H^{1/2}},\qquad
 \norm{\nabla v}_3+\norm{(-\Delta)^{1/2}v}_3
       \le C\norm{v}_{\dot H^{3/2}}.
\end{equation}
We also use $H^1\hookrightarrow L^6$ and $H^2\hookrightarrow L^\infty$.
\end{lemma}
\begin{proof}
Write $\langle\ell\rangle=(1+4\pi^2|\ell|^2)^{1/2}$.
The triangle inequality in $\R^4$ and the elementary power inequality give
\[
 \langle\ell\rangle^r
 \le C_r\bigl(\langle j\rangle^r+\langle\ell-j\rangle^r\bigr)
 \qquad(r\ge0).
\]
Using the latter inequality with $r=k$ in the Fourier convolution of $zw$ yields
\[
 \norm{zw}_{H^k}\le C_k\bigl(
   \norm{\widehat z}_{\ell^1}\norm{w}_{H^k}
  +\norm{\widehat w}_{\ell^1}\norm{z}_{H^k}\bigr).
\]
Here the sequence convolution estimate follows directly by writing an
$\ell^1*\ell^2$ convolution as a sum of translates and applying the triangle
inequality in $\ell^2$. Cauchy--Schwarz gives
\[
 \norm{\widehat z}_{\ell^1}
 \le\Bigl(\sum_{\ell\in\mathbb Z^3}\langle\ell\rangle^{-4}\Bigr)^{1/2}
      \norm{z}_{H^2}\le C\norm{z}_{H^2}.
\]
The sum converges: a dyadic shell of radius $2^j$ contains at most $C2^{3j}$
lattice points and its weighted contribution is at most $C2^{-j}$.
This proves \eqref{eq:algebra} and \eqref{eq:tame}, initially for Fourier
polynomials and then by their Fourier-norm approximation. The same absolute
Fourier convergence proves $\norm{z}_\infty\le C\norm{z}_{H^2}$.

Lemma~\ref{lem:critical-embeddings} in Appendix~\ref{sec:critical-appendix}
proves the critical embeddings on both domains. Apply its order-$1/2$ inequality to $v$,
each component of $\nabla v$, and $(-\Delta)^{1/2}v$ to obtain
\eqref{eq:embeddings}. Its order-one inequality applied to $z-\widehat z(0)$,
together with $|\widehat z(0)|\le\norm{z}_2$, gives the inhomogeneous
$H^1\hookrightarrow L^6$ assertion. Thus no integer
Gagliardo--Nirenberg estimate is needed here.
\end{proof}

\begin{proposition}[Local solutions, uniqueness and continuation]\label{prop:local}
Let $a\in\XX$ and let $f$ be periodic and smooth on every compact interval in $[0,\infty)$.
There is a unique maximal smooth solution of \eqref{eq:NS}, with normalized periodic pressure, on $[0,T_{\max})$.
If $S<\infty$ and
\begin{equation}\label{eq:criterion}
 \int_0^S\norm{u(t)}_{H^2}^2\dd t<\infty
\end{equation}
on a classical solution defined for $0\le t<S$, then that solution extends smoothly beyond $S$.
\end{proposition}
\begin{proof}
\emph{One common $H^3$ interval.}
Let $J=(1-\Delta)^{1/2}$, defined by the Fourier weights above. For $r\ge0$,
\[
 \norm{e^{\nu t\Delta}\PP\nabla\cdot A}_{H^r}
 \le C(\nu t)^{-1/2}\norm{A}_{H^r},\qquad t>0,
\]
because the multiplier has norm at most
$\sup_{\xi\in\R^3}|\xi|e^{-\nu t|\xi|^2}\le C(\nu t)^{-1/2}$.
The heat semigroup is a contraction and is strongly continuous on every
$H^r$, as follows by dominated convergence in its Fourier series.
Consider the mild equation
\begin{align}
 u(t)={}&e^{\nu t\Delta}a
 -\int_0^t e^{\nu(t-s)\Delta}\PP\nabla\cdot(u\otimes u)(s)\dd s \nonumber\\
 &+\int_0^t e^{\nu(t-s)\Delta}\PP f(s)\dd s.\label{eq:mild}
\end{align}
It maps the space of real, divergence-free fields in
$C([0,\tau];H^3)$ to itself. The integral is continuous in time: split off an
interval of length $h$ next to the upper limit, whose norm is $O(h^{1/2})$,
and use semigroup strong continuity on the remaining interval.
Lemma~\ref{lem:calculus} bounds the nonlinear integral by
$C\nu^{-1/2}\tau^{1/2}\norm{u}_{C H^3}^2$, and the difference of two such
integrals by
\[
 C\nu^{-1/2}\tau^{1/2}
 (\norm{u}_{C H^3}+\norm{v}_{C H^3})\norm{u-v}_{C H^3}.
\]
Take $R=2(\norm{a}_{H^3}+1)$ and choose $\tau>0$ so that
$\int_0^\tau\norm{f(s)}_{H^3}\dd s\le1$ and
$2C\nu^{-1/2}\tau^{1/2}R\le1/2$.
The map preserves the closed ball of radius $R$ and is a contraction there.
Its iterates are Cauchy by a geometric estimate in the complete space
$C([0,\tau];H^3)$; their limit is the unique fixed point in that ball.

\emph{Orthogonal truncations and uniform higher-order bounds.}
For an integer $N\ge1$ define
\[
 \widehat{Q_N z}(\ell)=\mathbf1_{\{|\ell|\le N\}}\widehat z(\ell).
\]
This is a self-adjoint orthogonal projection, has norm at most one on every
$H^r$, and commutes with $J$, $\PP$, spatial derivatives and the heat
semigroup. Its symmetric frequency set preserves real fields.
On exactly the same $C H^3$ ball and interval, the projected mild map gives
$u_N=Q_Nu_N$ solving
\[
 \partial_tu_N-\nu\Delta u_N
 =-Q_N\PP\nabla\cdot(u_N\otimes u_N)+Q_N\PP f,
 \qquad u_N(0)=Q_Na,
 \qquad\norm{u_N}_{C H^3}\le R.
\]
Differentiating the finitely many mild coefficient equations first gives
$C^1$ coefficients solving a polynomial ordinary differential equation;
repeated differentiation, since the forcing is smooth, gives $C^\infty$ coefficients. Thus $u_N$ is smooth and the following energy
identities are justified directly. For an integer $k\ge3$
put $X_{k,N}=\norm{u_N}_{H^k}$ and
$D_{k,N}=\norm{\nabla u_N}_{H^k}$. Apply $J^k$, pair with $J^ku_N$, and
move both projections to $u_N$ using self-adjointness. Integrating the
nonlinear divergence by parts and using \eqref{eq:tame} gives
\[
 \tfrac12(X_{k,N}^2)'+\nu D_{k,N}^2
 \le C_k\norm{u_N}_{H^2}X_{k,N}D_{k,N}
       +\norm{f}_{H^k}X_{k,N}.
\]
The dissipation is the displayed gradient norm; no Poincar\'e estimate on the
possibly nonzero mean is used. Young's inequality and division regularized
by $(X_{k,N}^2+\zeta^2)^{1/2}$ imply, after integration and
$\zeta\downarrow0$,
\[
 X_{k,N}(t)\le
 e^{C_{k,\nu}R^2\tau}
 \left(\norm{a}_{H^k}+\int_0^\tau\norm{f(s)}_{H^k}\dd s\right)=:M_k.
\]
Every $M_k$ is independent of $N$. In particular, higher regularity does not
require shrinking the original $H^3$ interval.

\emph{Strong convergence, without a compact embedding.}
Subtract the mild equation for $u$ from the projected equation. With the
same contraction constant $q\le1/2$, the difference estimate is
\begin{align*}
 \norm{u_N-u}_{C H^3}\le{}&q\norm{u_N-u}_{C H^3}
  +\norm{(I-Q_N)a}_{H^3}
  +\int_0^\tau\norm{(I-Q_N)\PP f(s)}_{H^3}\dd s\\
 &+C\nu^{-1/2}\tau^{1/2}
       \norm{(I-Q_N)(u\otimes u)}_{C H^3}.
\end{align*}
The last three terms tend to zero. For the last one, $u\otimes u$ is
continuous into $H^3$, so its image of $[0,\tau]$ is compact. Strong
convergence $Q_N\to I$ is uniform on that image: cover it by finitely many
balls of radius $\varepsilon$, use convergence at their centers, and use
$\norm{I-Q_N}\le1$. The forcing term follows either in the same way or by
dominated convergence. Hence $u_N\to u$ in $C H^3$.

At each time, convergence of the Fourier coefficients and Fatou's lemma
transfer $\norm{u_N(t)}_{H^k}\le M_k$ to
$\norm{u(t)}_{H^k}\le M_k$. For $3<m<k$, H\"older's inequality in the
Fourier sum gives the explicit interpolation bound
\[
 \norm{w}_{H^m}
 \le\norm{w}_{H^3}^{(k-m)/(k-3)}
       \norm{w}_{H^k}^{(m-3)/(k-3)}.
\]
Indeed each squared weighted Fourier term is the product of the
order-three and order-$k$ terms raised to those two exponents; summing and
applying H\"older proves the formula. Apply it to $w=u_N-u$, whose $H^k$
norm is at most $2M_k$, and choose an integer $k>m$. We obtain convergence in
$C([0,\tau];H^m)$ for every fixed integer $m\ge0$ (orders at most three
follow directly). Thus $u$ has continuous spatial regularity of every order
on the same interval.

\emph{Time regularity and pressure.}
The right sides of the truncated evolution equations converge in
$C H^m$ for every $m$, using the just-proved convergence in $C H^{m+3}$,
the product estimate and uniform projection convergence on continuous
images. Passing to their time integrals gives
\[
 u\in C^1([0,\tau];H^m),\qquad
 \partial_tu=\nu\Delta u-\PP\nabla\cdot(u\otimes u)+\PP f.
\]
Time derivatives at zero are right derivatives. Repeatedly differentiate
this identity: the derivative of the product is a finite sum of products
of previously obtained derivatives. Induction, with arbitrarily high
spatial order available at each step, gives
$u\in C^j([0,\tau];H^m)$ for all $j,m$.
The pressure formula preceding the lemma has, at nonzero frequencies,
order $-1$ on $f$ and order zero on $u\otimes u$; its zero coefficient is
set to zero. It therefore gives the same time and spatial regularity for
$p$. Finally, for $m>r+3/2$, Cauchy--Schwarz makes the Fourier series of
all spatial derivatives of order at most $r$ absolutely and uniformly
convergent, since
$\sum_\ell\langle\ell\rangle^{-2(m-r)}<\infty$.
Applied also to time derivatives, this proves joint classical smoothness
on every such closed interval.

\emph{Uniqueness and the maximal interval.}
If $u_1,u_2$ are two smooth solutions with the same data and $z=u_1-u_2$,
then
\[
 \partial_tz-\nu\Delta z+(u_1\cdot\nabla)z
 +(z\cdot\nabla)u_2+\nabla(p_1-p_2)=0.
\]
Testing against $z$ and using incompressibility and periodicity gives
\[
 \tfrac12(\norm{z}_2^2)'+\nu\norm{\nabla z}_2^2
 \le\norm{\nabla u_2}_\infty\norm{z}_2^2.
\]
Gr\"onwall yields $z=0$ on every closed common interval, and the pressure
normalization identifies the pressures. Restarting the construction at
any smooth time and joining these unique solutions defines the maximal
interval $[0,T_{\max})$.

\emph{Continuation.}
The same energy calculation on a smooth solution, now writing
$X_k=\norm{u}_{H^k}$ and $D_k=\norm{\nabla u}_{H^k}$, gives
\[
 \tfrac12(X_k^2)'+\nu D_k^2
 \le C_k\norm{u}_{H^2}X_kD_k+\norm{f}_{H^k}X_k.
\]
Young's inequality and regularized division therefore give the integrated
form of
\begin{equation}\label{eq:highcontinuation}
 X_k'\le C_{k,\nu}\norm{u}_{H^2}^2X_k+\norm{f}_{H^k}.
\end{equation}
Under \eqref{eq:criterion}, Gr\"onwall bounds $X_k$ uniformly for $t<S$ at
every fixed $k$, since $f$ is smooth on $[0,S+1]$.
The $H^3$ construction gives a positive lower bound for the restart interval
using only a bound for $\norm{u(t_0)}_{H^3}$ and
$\sup_{[0,S+1]}\norm{f}_{H^3}$. Choose $t_0<S$ closer to $S$ than that
interval length (and shorten the length to at most one). The restarted
solution extends beyond $S$, has all higher regularity by the common
interval argument, and equals the original solution on their overlap by
uniqueness. This proves continuation without assuming an endpoint value
in $H^3$. The criterion proved here is precisely \eqref{eq:criterion};
the unbounded speed of the inserted solutions below follows separately
from the packet theorem.
\end{proof}

\section{The compact packet and its scaling}\label{sec:packet}
\begin{theorem}[Compact forced blowup packet, OpenAI]\label{thm:packet}
At the fixed viscosity $\nu>0$, there exist smooth fields $U,P$ on
$\R^3\times[0,1)$, a force
$F\in C_c^\infty(\R^3\times(0,\infty);\R^3)$, and a compact set $K$ such that
\begin{align}
 &\partial_tU+(U\cdot\nabla)U-\nu\Delta U+\nabla P=F,\qquad
 \nabla\cdot U=0,\qquad U(0)=0,\label{eq:packet}\\
 &\supp U(t)\cup\supp P(t)\subset K\quad(0\le t<1),\nonumber\\
 &\sup_{0\le t<1}\norm{U(t)}_2<\infty,\qquad
 \limsup_{t\uparrow1}\norm{U(t)}_\infty=\infty.\label{eq:packetblowup}
\end{align}
\end{theorem}
This is \cite[Theorem 1.1]{openai}. Its associated formalization \cite{lean} supplies the zero-data witness and compact-support properties. We use precisely the displayed packet properties to derive the estimates and distribution results below.

\begin{lemma}[Energy, dissipation and initial vanishing]\label{lem:packetenergy}
The packet satisfies
\[
 M:=\sup_{0\le t<1}\norm{U(t)}_2<\infty,\qquad
 D:=\norm{\nabla U}_{L^2((0,1)\times\R^3)}<\infty.
\]
Both $U$ and $P$ vanish on an initial time interval and therefore extend smoothly by zero to negative times.
\end{lemma}
\begin{proof}
On each closed interval $[0,b]$ with $b<1$, compact spatial support justifies the energy identity
\[
 \tfrac12(\norm{U(t)}_2^2)'
 +\nu\norm{\nabla U(t)}_2^2=\ip{F(t)}{U(t)}.
\]
Set $N(t)=\int_0^t\norm{F(s)}_2\dd s$. Dropping dissipation gives
$\tfrac12(\norm{U}_2^2)'\le\norm{F}_2\norm{U}_2$.
For $\zeta>0$, the derivative of $(\norm{U}_2^2+\zeta^2)^{1/2}$ is at most $\norm{F}_2$. Integrating from zero and letting $\zeta\downarrow0$ proves $\norm{U(t)}_2\le N(t)$.

Substituting this bound into the integrated energy identity gives
\begin{equation}\label{eq:packetenergy}
 \norm{U(t)}_2^2+2\nu\int_0^t\norm{\nabla U(s)}_2^2\dd s
 \le 2\int_0^t\norm{F(s)}_2N(s)\dd s=N(t)^2.
\end{equation}
Since $F$ is smooth and time compact, the right side is bounded independently of $t<1$. Monotone convergence of the dissipation integral proves $D<\infty$.

The temporal support of $F$ is a compact subset of $(0,\infty)$, so $F=0$ on $[0,\tau]$ for some $\tau>0$. Estimate~\eqref{eq:packetenergy} gives $U=0$ there. Equation~\eqref{eq:packet} then implies $\nabla P=0$, and compact spatial support fixes this spatial constant to zero. Vanishing on an interval proves smoothness of both zero extensions.
\end{proof}

\begin{lemma}[Uniform localization of fractional norms]\label{lem:localization}
Let $B$ be a fixed coordinate ball whose closure lies inside a torus fundamental cube. If $z$ is smooth and supported in $B$, let $z_{\R}$ be its zero extension in that cube to $\R^3$, and let $z_{\TT}$ be its periodization. For $0<s<1$,
\begin{equation}\label{eq:localization}
 \norm{z_{\TT}}_{H^s(\TT)}
 \le C_{s,B}\bigl(\norm{z_{\R}}_2+
                   \norm{z_{\R}}_{\dot H^s(\R^3)}\bigr).
\end{equation}
The constant is uniform as the support shrinks inside $B$. At $s=0$ the $L^2$ norms are equal; at $s=1$ the corresponding equality holds for the $L^2$ gradient norms.
\end{lemma}
\begin{proof}
We first identify the difference expressions with the Fourier norms used
in this article. Define
\[
 c_s=\int_{\R^3}\frac{|e^{ih_1}-1|^2}{|h|^{3+2s}}\dd h.
\]
It is positive and finite: near zero the numerator is at most $|h|^2$,
giving a radial integral $\int_0^1r^{1-2s}\dd r$; away from zero it is at
most four, giving $\int_1^\infty r^{-1-2s}\dd r$.
For a compact smooth $Z$ on $\R^3$, change variables $y=x+h$, then use
Plancherel and Tonelli on the nonnegative integrands to obtain
\begin{align*}
 I_{\R}(Z)&:=\int_{\R^3}\int_{\R^3}
       \frac{|Z(x+h)-Z(x)|^2}{|h|^{3+2s}}\dd x\dd h\\
 &=\int_{\R^3}|\widehat Z(\xi)|^2
       \left(\int_{\R^3}\frac{|e^{i\xi\cdot h}-1|^2}
                                      {|h|^{3+2s}}\dd h\right)\dd\xi
 =c_s\norm{Z}_{\dot H^s(\R^3)}^2.
\end{align*}
For $\xi\ne0$, rotation and dilation identify the inner integral with
$c_s|\xi|^{2s}$; it is zero when $\xi=0$.

Let $Q$ be the chosen unit fundamental cube and define the nonnegative
periodic kernel, off its lattice singularities, by
\[
 K_s(h)=\sum_{n\in\mathbb Z^3}|h+n|^{-3-2s}.
\]
For a smooth periodic $z$, partitioning $\R^3$ into translates of $Q$ and
using periodicity gives
\begin{align*}
 I_{\TT}(z)&:=\int_Q\int_Q|z(x)-z(y)|^2K_s(x-y)\dd x\dd y\\
 &=\int_{\R^3}\frac{\norm{z(\,\cdot+h)-z}_{L^2(Q)}^2}
                            {|h|^{3+2s}}\dd h\\
 &=\sum_{\ell\in\mathbb Z^3}|\widehat z(\ell)|^2
       \int_{\R^3}\frac{|e^{2\pi i\ell\cdot h}-1|^2}
                            {|h|^{3+2s}}\dd h
 =c_s\norm{z-\widehat z(0)}_{\dot H^s(\TT)}^2.
\end{align*}
The second line uses ordinary periodic $L^2$ norms. Parseval proves the
third line and Tonelli justifies interchanging the nonnegative sum and
integral. Thus both identifications have the same explicitly specified
constant $c_s$, with the frequency factor $2\pi$ retained on the torus.
Also $(1+r^2)^s\le1+r^{2s}$ for $0<s<1$, so
\[
 \norm{z}_{H^s(\TT)}^2
 \le\norm{z}_2^2+c_s^{-1}I_{\TT}(z).
\]

It remains to compare the two difference integrals uniformly. Set
$d=\operatorname{dist}(\overline B,\partial Q)>0$.
For $x\in B$, $y\in Q$ and $n\ne0$ we have
$|x-y+n|\ge d$: the point $y-n$ lies outside $Q$ up to its boundary.
For $|n|>2\sqrt3$ the same distance is at least $|n|/2$.
Consequently
\[
 \sup_{x\in B,\,y\in Q}
 \sum_{n\ne0}|x-y+n|^{-3-2s}\le C_{s,d}<\infty,
\]
and the identical estimate holds with $x,y$ reversed. The $n=0$ term of
$I_{\TT}(z_{\TT})$ is at most $I_{\R}(z_{\R})$.
For all the other terms use
$|z(x)-z(y)|^2\le2|z(x)|^2+2|z(y)|^2$; each nonzero term on the right has
its corresponding point in $B$. As $|Q|=1$, their total is at most
$4C_{s,d}\norm{z_{\R}}_2^2$.
This proves \eqref{eq:localization}. Both $d$ and all constants refer to
$B,Q,s$, independently of any smaller support. At orders zero and one,
integration over the single copy of the support proves the stated $L^2$
and gradient identities.
\end{proof}

Choose a compact set $K_*$ containing $K$ and the spatial projection of $\supp F$. Fix a point $x_0$ in a torus coordinate ball. For sufficiently small $\eps>0$, require
\[
 2\eps^2<T,\qquad x_0+\eps K_*\subset B,
 \qquad t_\eps=T-\eps^2,
\]
and define, using the zero extensions from Lemma~\ref{lem:packetenergy},
\begin{align}
 U_\eps(x,t)&=\eps^{-1}U\left(\frac{x-x_0}{\eps},
                              \frac{t-t_\eps}{\eps^2}\right),\nonumber\\
 P_\eps(x,t)&=\eps^{-2}P\left(\frac{x-x_0}{\eps},
                              \frac{t-t_\eps}{\eps^2}\right),\label{eq:scaling}\\
 F_\eps(x,t)&=\eps^{-3}F\left(\frac{x-x_0}{\eps},
                              \frac{t-t_\eps}{\eps^2}\right).\nonumber
\end{align}
The first two fields are used for $t<T$; the force is defined for all positive times. Place this single whole-space packet in $B$ and then periodize. This operation does not rescale an already periodic field, which would replicate its support.

\begin{proposition}[Packet scaling at fixed viscosity]\label{prop:scaling}
The fields in \eqref{eq:scaling} solve the periodic momentum equation at viscosity $\nu$, start from zero, and have unbounded speed at $T$. A spatially constant pressure adjustment enforces the mean-zero normalization. Their norms satisfy
\begin{align}
 \norm{U_\eps}_{L^\infty(0,T;L^2)}
    &=\eps^{1/2}M,&
 \norm{\nabla U_\eps}_{L^2(0,T;L^2)}
    &=\eps^{1/2}D,\label{eq:packetEscale}\\
 \norm{F_\eps}_{L^q(0,\infty;L^p)}
    &=\eps^{\alpha(p,q)}\norm{F}_{L^q(0,\infty;L^p)},&
 \alpha(p,q)&=-3+\frac3p+\frac2q,\label{eq:packetFscale}
\end{align}
for $1\le p,q\le\infty$. For $0\le s\le1$,
\begin{equation}\label{eq:packetHs}
 \norm{F_\eps}_{L^1(0,\infty;H^s(\TT))}
 \le C_s\bigl(\eps^{1/2}+\eps^{1/2-s}\bigr).
\end{equation}
In particular the force tends to zero in $L^1_tH^s_x$ for every $s<1/2$, including negative $s$.
\end{proposition}
\begin{proof}
Writing $y=(x-x_0)/\eps$ and $\sigma=(t-t_\eps)/\eps^2$, every term of the momentum equation gains the common factor $\eps^{-3}$. Incompressibility is preserved and no viscosity rescaling occurs. The temporal extension is smooth by initial vanishing. The speed identity
$\norm{U_\eps(t)}_\infty=\eps^{-1}\norm{U(\sigma)}_\infty$
proves unboundedness as $t\uparrow T$.

At a fixed time the velocity has amplitude factor $\eps^{-1}$ and spatial volume factor $\eps^3$, giving the first identity in \eqref{eq:packetEscale}. Its gradient has amplitude factor $\eps^{-2}$. Squaring and integrating in both space and time gives
$\eps^{-4}\eps^3\eps^2=\eps$ times the original dissipation integral, proving the second identity.

For finite $p,q$, the force's spatial norm gains $\eps^{-3+3/p}$ and its time norm gains $\eps^{2/q}$. The positive-time domain includes the whole transformed support because $t_\eps>0$ and the original force vanishes at negative times. This proves \eqref{eq:packetFscale}; the essential-supremum cases follow by the same change of variables with zero reciprocal exponents.

Euclidean homogeneous Sobolev scaling gives, for $0<s<1$,
\[
 \norm{\eps^{-3}F((\,\cdot-x_0)/\eps,\sigma)}
          _{\dot H^s(\R^3)}
 =\eps^{-3/2-s}\norm{F(\cdot,\sigma)}_{\dot H^s(\R^3)}.
\]
One may verify the exponent either by Fourier transform or by the double integral in Lemma~\ref{lem:localization}: its squared scaling factor is
$\eps^{-6}\eps^6\eps^{-3-2s}=\eps^{-3-2s}$.
Time integration contributes $\eps^2$. The inhomogeneous $L^2$ term has order $\eps^{1/2}$. Lemma~\ref{lem:localization} therefore proves \eqref{eq:packetHs}. At $s=0$ use the mixed-norm identity, and at $s=1$ use the gradient scaling. For $s<0$, the Fourier definition gives $\norm{z}_{H^s}\le\norm{z}_2$, so the $s=0$ estimate proves convergence without any homogeneous negative-order claim.
\end{proof}

The force peak grows as $\eps^{-3}\norm{F}_\infty$. Thus the smallness obtained here is smallness of an integrated norm. At the critical exponent, the Euclidean $L^1_t\dot H^{1/2}_x$ norm is scale invariant; the independent regularity argument in Section~\ref{sec:critical} will turn that scaling boundary into a genuine non-density statement.

\section{Exact insertion into a regular background}\label{sec:insertion}
Fix a reference solution $(v,\pi,g)$ smooth on $[0,T+\delta]$, where
$\delta>0$, $g\in\FF$, and $v(0)=a\in\XX$. Direct addition of $U_\eps$ would produce cross terms involving the singular packet and the nonzero reference. Those terms need not extend smoothly through $T$. We instead remove the reference in a small neighborhood of the packet support before adding the packet.

\begin{lemma}[A local divergence-free cutoff]\label{lem:potential}
Let $v$ be smooth and divergence free in a spatial ball centered at $x_0$. In that ball, with $y=x-x_0$, define
\begin{equation}\label{eq:potential}
 A(x,t)=\int_0^1 r\,v(x_0+ry,t)\times y\dd r.
\end{equation}
Then $\nabla\times A=v$. Choose $\theta\in C_c^\infty(\R^3)$ equal to one on a neighborhood of $K_*$, and choose
$\eta\in C_c^\infty((-2,2))$ equal to one on $[-1,1]$. Put
\begin{equation}\label{eq:cutoff}
 \theta_\eps(x)=\theta((x-x_0)/\eps),\quad
 \eta_\eps(t)=\eta((t-T)/\eps^2),\quad
 w_\eps=-\nabla\times(\eta_\eps\theta_\eps A).
\end{equation}
For sufficiently small $\eps$, $w_\eps$ is a smooth, divergence-free field supported inside the coordinate ball and in
$(T-2\eps^2,T+2\eps^2)$. During the active presingular interval of $U_\eps$,
\begin{equation}\label{eq:bgzero}
 v+w_\eps=0
 \quad\hbox{on an open neighborhood of }\supp U_\eps(t).
\end{equation}
\end{lemma}
\begin{proof}
Use the vector identity
\[
 \nabla\times(V\times y)
 =V\,\nabla\cdot y-y\,\nabla\cdot V
   +(y\cdot\nabla)V-(V\cdot\nabla)y.
\]
For $V(y)=v(x_0+ry,t)$, incompressibility gives $\nabla_y\cdot V=0$,
$\nabla_y\cdot y=3$, and $(V\cdot\nabla_y)y=V$.
Consequently,
\begin{align*}
 \nabla_y\times\bigl(r\,v(x_0+ry,t)\times y\bigr)
 &=2r\,v(x_0+ry,t)+r^2(y\cdot\nabla_x)v(x_0+ry,t)\\
 &=\frac{\dd}{\dd r}\bigl(r^2v(x_0+ry,t)\bigr).
\end{align*}
Integrating from $r=0$ to $r=1$ proves $\nabla\times A=v$.

Take $\eps$ small enough that the scaled support of $\theta$ is strictly inside the coordinate ball and $2\eps^2<\min(T,\delta)$. Since the spatial cutoff is zero near the boundary of the coordinate ball, extending \eqref{eq:cutoff} by zero gives a globally smooth periodic field. Its divergence vanishes because it is a curl. The temporal support follows directly from $\eta_\eps$.

The packet is zero before $t_\eps=T-\eps^2$. For
$t_\eps\le t<T$, the variable $(t-T)/\eps^2$ belongs to $[-1,0)$, on which $\eta=1$. On a neighborhood of the packet support, $\theta_\eps=1$ and its derivatives vanish. There $w_\eps=-\nabla\times A=-v$, proving \eqref{eq:bgzero}.
\end{proof}

\begin{lemma}[Bounds for the background correction]\label{lem:correction}
For $w_\eps$ in Lemma~\ref{lem:potential}, define
\begin{align}
 H_\eps={}&\partial_tw_\eps-\nu\Delta w_\eps
 +(v\cdot\nabla)w_\eps+(w_\eps\cdot\nabla)v
 +(w_\eps\cdot\nabla)w_\eps.\label{eq:H}
\end{align}
The force $H_\eps$ is smooth across $T$ and extends by zero to all times outside its cutoff support. Its spatial support has volume $O(\eps^3)$, and its temporal support has length $O(\eps^2)$. For every spatial multi-index $\beta$ and integer $j\ge0$,
\begin{equation}\label{eq:derivativebounds}
 |\partial_t^j\partial_x^\beta w_\eps|
 \le C_{\beta,j}\eps^{-2j-|\beta|},\qquad
 |\partial_x^\beta H_\eps|
 \le C_\beta\eps^{-2-|\beta|}.
\end{equation}
In particular,
\begin{align}
 \norm{w_\eps}_{E_T}&\le C\eps^{3/2},\label{eq:wE}\\
 \norm{H_\eps}_{L^q(0,\infty;L^p)}
   &\le C_{p,q}\eps^{-2+3/p+2/q}
     =C_{p,q}\eps^{\alpha(p,q)+1},\label{eq:Hmixed}\\
 \norm{H_\eps}_{L^1(0,\infty;H^s)}
   &\le C_s\bigl(\eps^{3/2}+\eps^{3/2-s}\bigr)
       \quad(0\le s\le1).\label{eq:HHs}
\end{align}
All constants are independent of sufficiently small $\eps$.
\end{lemma}
\begin{proof}
Introduce $z=(x-x_0)/\eps$ and $\sigma=(t-T)/\eps^2$. Then
\[
 A(x_0+\eps z,T+\eps^2\sigma)=\eps\,\mathcal A_\eps(z,\sigma),
\]
where
\[
 \mathcal A_\eps(z,\sigma)
 =\int_0^1r\,v(x_0+\eps rz,T+\eps^2\sigma)\times z\dd r.
\]
On the fixed support of $\theta(z)\eta(\sigma)$, every fixed-order derivative of $\mathcal A_\eps$ is bounded uniformly in $\eps$. Differentiation of $v$ in $z$ or $\sigma$ introduces nonnegative powers of $\eps$, and the reference is smooth on a fixed neighborhood of the relevant spacetime set. Thus
\[
 w_\eps(x_0+\eps z,T+\eps^2\sigma)
 =W_\eps(z,\sigma)
 :=-\nabla_z\times\bigl(\eta(\sigma)\theta(z)\mathcal A_\eps(z,\sigma)\bigr)
\]
is a uniformly smooth family supported in a fixed cylinder. This proves the first bound in \eqref{eq:derivativebounds}.

To check the force bound, set
$V_\eps(z,\sigma)=v(x_0+\eps z,T+\eps^2\sigma)$.
Formula~\eqref{eq:H} becomes
\begin{align*}
 H_\eps(x_0+\eps z,T+\eps^2\sigma)
 =\eps^{-2}\bigl[&
 \partial_\sigma W_\eps-\nu\Delta_zW_\eps
 +\eps(V_\eps\cdot\nabla_z)W_\eps\\
 &+\eps^2(W_\eps\cdot\nabla_x)v(x_0+\eps z,T+\eps^2\sigma)
 +\eps(W_\eps\cdot\nabla_z)W_\eps\bigr].
\end{align*}
The expression in brackets and all its fixed spatial derivatives are uniformly bounded on the fixed cylinder. This proves the second bound in \eqref{eq:derivativebounds}. It also makes explicit that $H_\eps$ involves only the smooth reference and fixed cutoffs, not derivatives of a singular field.

From the support volume and duration,
\[
 \norm{w_\eps}_{L^\infty_tL^2_x}\le C\eps^{3/2},\qquad
 \norm{\nabla w_\eps}_{L^2_tL^2_x}
 \le C\eps^{-1}\eps^{3/2}\eps=C\eps^{3/2}.
\]
This proves \eqref{eq:wE}. The amplitude bound $C\eps^{-2}$ gives
\eqref{eq:Hmixed} by the same support calculation, including the essential-supremum endpoints.

For \eqref{eq:HHs}, the rescaled force profiles have uniformly bounded spatial $H^s$ norms, supported on a fixed spatial set during a fixed interval of $\sigma$. An amplitude $\eps^{-2}$ gives the spatial homogeneous factor $\eps^{-1/2-s}$; time integration adds $\eps^2$. Thus the homogeneous contribution is $O(\eps^{3/2-s})$, while the $L^2$ contribution is $O(\eps^{3/2})$. Lemma~\ref{lem:localization} transfers these estimates to the torus. Smooth zero extension in time is valid because $\eta$, with all its derivatives, vanishes near the fixed interval endpoints $\sigma=-2$ and $\sigma=2$, where the reference remains smooth.
\end{proof}

\Needspace{14\baselineskip}
\begin{theorem}[Exact local insertion]\label{thm:insertion}
Let $(v,\pi,g)$ be a periodic solution smooth on $[0,T+\delta]$, with
$\delta>0$, $g\in\FF$ and $v(0)=a\in\XX$. Fix any nonempty coordinate ball.
For all sufficiently small $\eps>0$, there are $g_\eps\in\FF$ and a solution $u_\eps$ with
\begin{enumerate}[label=(\roman*),nosep]
\item $T_{\max}^{\nu}(a,g_\eps)=T$ and
$\limsup_{t\uparrow T}\norm{u_\eps(t)}_\infty=\infty$;
\item $u_\eps=v$ for $0\le t\le T-2\eps^2$;
\item $u_\eps-v$ divergence free and supported, for every $t<T$, in a ball of diameter $O(\eps)$ inside the chosen ball;
\item the simultaneous bounds
\end{enumerate}
\begin{align}
 \norm{u_\eps-v}_{E_T}
   &\le (M+D)\eps^{1/2}+C\eps^{3/2},\label{eq:Eclose}\\
 \norm{g_\eps-g}_{L^q_tL^p_x}
   &\le C_{p,q}\bigl(\eps^{\alpha(p,q)}
                    +\eps^{\alpha(p,q)+1}\bigr),\label{eq:Fclose}\\
 \norm{g_\eps-g}_{L^1_tH^s_x}
   &\le C_s\bigl(\eps^{1/2-s}+\eps^{3/2-s}\bigr)
                         \quad(0\le s<1/2).\label{eq:Hsclose}
\end{align}
Here force norms are taken over $(0,\infty)$, and
$\alpha(p,q)=-3+3/p+2/q$ for $1\le p,q\le\infty$.
For $s<0$, the force difference also tends to zero in $L^1_tH^s_x$.
\end{theorem}
\begin{proof}
Choose the scaled packet and the correction from Lemmas~\ref{lem:potential}--\ref{lem:correction}, with all supports inside the prescribed ball, and set
\begin{equation}\label{eq:insertion}
 u_\eps=v+w_\eps+U_\eps,\qquad
 p_\eps=\pi+P_\eps,\qquad
 g_\eps=g+H_\eps+F_\eps.
\end{equation}
Initially use the compact representative $P_\eps$; subtract its spatial mean to normalize $p_\eps$.

Let $b_\eps=v+w_\eps$. Expanding the reference equation and \eqref{eq:H} gives
\[
 \partial_tb_\eps+(b_\eps\cdot\nabla)b_\eps
 -\nu\Delta b_\eps+\nabla\pi=g+H_\eps.
\]
Adding the packet equation produces the equation for $u_\eps$ except for
\[
 (b_\eps\cdot\nabla)U_\eps+(U_\eps\cdot\nabla)b_\eps.
\]
Both terms are identically zero. On a neighborhood of $\supp U_\eps$ this follows from $b_\eps=0$, including its derivatives, by \eqref{eq:bgzero}. Outside that support the packet and its derivatives vanish; smoothness ensures the same assertion at the support boundary. When the packet is inactive it is zero everywhere. Thus \eqref{eq:NS} holds exactly for \eqref{eq:insertion} on $[0,T)$.

Each summand of the velocity is divergence free. The correction is zero for
$t\le T-2\eps^2$, and the packet is zero before $T-\eps^2$, so the initial datum and the stated earlier history are unchanged. On the active packet support, $u_\eps=U_\eps$. Hence
\[
 \norm{u_\eps(t)}_\infty\ge\norm{U_\eps(t)}_\infty
\]
and the maximum norm is unbounded as $t\uparrow T$.

The velocity is smooth on every closed presingular interval. The force $F_\eps$ is globally smooth by its original definition, and $H_\eps$ is globally smooth by Lemma~\ref{lem:correction}. Their time supports remain compact subsets of $(0,\infty)$ when $\eps$ is sufficiently small. Therefore $g_\eps\in\FF$, independently of the absence of a classical extension of $u_\eps$ at $T$. Uniqueness in Proposition~\ref{prop:local} identifies the constructed velocity with the maximal solution for $(a,g_\eps)$ up to every $T'<T$. Its lifespan is at least $T$, and the unbounded maximum norm excludes extension beyond $T$. Thus the lifespan is exactly $T$.

Finally, the triangle inequality, Proposition~\ref{prop:scaling} and
Lemma~\ref{lem:correction} give \eqref{eq:Eclose}--\eqref{eq:Hsclose}. In
\eqref{eq:Hsclose}, the lower-order inhomogeneous terms are absorbed because
$0<\eps\le1$ and $s\ge0$. For negative $s$, use
$\norm{z}_{H^s}\le\norm{z}_2$ and the $s=0$ estimate. This proves every assertion.
\end{proof}

\begin{corollary}[Interior no-slip insertion]\label{cor:boundary}
Let $\Omega\subset\R^3$ be a bounded box or a bounded smooth domain. Suppose
$(v,\pi,g)$ is a smooth no-slip solution through $T+\delta$, and $g$ is smooth on $\overline\Omega\times[0,\infty)$ with temporal support compact in $(0,\infty)$.
Here smoothness on a closed spacetime slab means restriction of a
$C^\infty$ field from an open neighborhood of that slab; this convention
also specifies smoothness at the edges and corners of a box. We assume
$v,\pi$ have this regularity on
$\overline\Omega\times[0,T+\delta]$, and $g$ on each finite closed slab.
The equation and incompressibility hold in $\Omega$, and
$v|_{\partial\Omega}=0$. Pressure may be normalized by zero spatial mean.
Existence of this compatible reference is an assumption of the corollary.
The construction in Theorem~\ref{thm:insertion} applies inside any prescribed interior ball, preserving the initial velocity and no-slip boundary values.
The energy estimates use $L^2(\Omega)$ and the force estimates use
$L^1(0,\infty;H^s(\Omega))$, with the restriction norm
\eqref{eq:restriction-norm}. For every real $s$, the domain and zero-extended
force-difference norms are comparable by \eqref{eq:zero-extension}, with a
constant independent of $\eps$; the convergence assertion remains $s<1/2$.
\end{corollary}
\begin{proof}
The vector potential and all cutoff estimates are local. Choose the correction and packet supports in a smaller closed ball strictly inside $\Omega$.
The computation of the momentum equation in Theorem~\ref{thm:insertion} is unchanged. The new velocity equals the reference in a fixed boundary collar for all presingular times, so the no-slip boundary values are preserved. The force correction is supported in the same interior region and has the same globally smooth time extension.

For $0\le s<1/2$, the Euclidean scaling proof directly bounds the zero
extensions of the force differences in $L^1(0,\infty;H^s(\R^3))$.
For $s<0$, use $\norm{E_0(g_\eps-g)}_{H^s(\R^3)}
\le\norm{E_0(g_\eps-g)}_{L^2(\R^3)}$ at each time.
All these force-difference supports, including those after $T$, lie in one
fixed compact interior ball $K$ for every sufficiently small $\eps$.
Consequently \eqref{eq:zero-extension} gives the asserted domain estimates
and the scale-independent comparison after time integration. The energy estimates follow by integrating over that ball. Finally, classical uniqueness on a no-slip domain follows from the same difference-energy calculation as in Proposition~\ref{prop:local}; its boundary terms vanish. Thus the constructed solution is singular exactly at $T$.
\end{proof}

\begin{remark}[Concentration and force amplitudes]\label{rem:peaks}
The force convergence in Theorem~\ref{thm:insertion} is compatible with diverging pointwise amplitudes. Since the original force is nonzero,
\[
 \norm{g_\eps-g}_{L^\infty_{t,x}}
 \ge\eps^{-3}\norm{F}_\infty-C\eps^{-2}\longrightarrow\infty.
\]
The first term is the exact packet peak; the second bounds the entire background-force correction. The original force must be nonzero, since \eqref{eq:packetenergy} would otherwise force $U=0$. Thus this lower bound is nonvacuous. Each $g_\eps$ is smooth, but this family has no uniform force-amplitude or derivative bounds.
\end{remark}

\section{Density, trajectory closure and data projections}\label{sec:density}
\begin{proposition}[Density on every fixed initial-velocity slice]\label{prop:density}
For $a\in\XX$, $\nu>0$, $T>0$ and $s<1/2$,
\[
 \forall g\in\FF\ \forall\rho>0\ \exists f\in\FF:
 \quad \norm{f-g}_{L^1_tH^s_x}<\rho,\qquad
       T_{\max}^{\nu}(a,f)\le T.
\]
\end{proposition}
\begin{proof}
Fix $g$ and $\rho$. There are two exhaustive cases.

If $T_{\max}^{\nu}(a,g)\le T$, choose $f=g$. It already belongs to
$\BB_{\nu,a,T}$ and the norm difference is zero.

If $T_{\max}^{\nu}(a,g)>T$, choose $\delta>0$ so that
$T+\delta<T_{\max}^{\nu}(a,g)$, and use its smooth solution as the reference in Theorem~\ref{thm:insertion}. For every sufficiently small $\eps$, that theorem supplies $g_\eps\in\BB_{\nu,a,T}$ with lifespan exactly $T$.
For $0\le s<1/2$, \eqref{eq:Hsclose} tends to zero. For $s<0$, its $s=0$ case and the embedding $L^2\hookrightarrow H^s$ do so. Choose $\eps$ so that the force difference is less than $\rho$ and set $f=g_\eps$.

The initial datum $a$ and the parameters $\nu,T$ are unchanged in both cases. This is precisely density in the relative force norm topology.
\end{proof}

\begin{corollary}[A sufficient mixed-norm region]\label{cor:mixed}
For every fixed $a\in\XX$, the set $\BB_{\nu,a,T}$ is dense in $\FF$ with its relative $L^q(0,\infty;L^p)$ topology whenever
\begin{equation}\label{eq:mixedregion}
 1\le p,q\le\infty,\qquad \frac3p+\frac2q>3.
\end{equation}
\end{corollary}
\begin{proof}
Use the same two cases as in Proposition~\ref{prop:density}. In the regular-reference case, \eqref{eq:Fclose} tends to zero because
$\alpha(p,q)>0$ and hence $\alpha(p,q)+1>0$.
\end{proof}
This region includes $L^1_tL^2_x$ and $L^2_tL^{4/3}_x$. It is a sufficient condition, not a classification of all mixed Lebesgue topologies. A nonpositive value of the packet's scaling exponent alone proves neither density nor non-density.

\begin{corollary}[Strong closure in energy and dissipation]\label{cor:closure}
Fix $a\in\XX$. Let $\mathcal R_{a,T}$ be the trajectories with initial velocity $a$ and forces in $\FF$ that are smooth through $T$.
Let $\mathcal S_{a,T}$ be the trajectories with the same initial velocity and forces in $\FF$ that are smooth on $[0,T)$, have finite $E_T$ norm, and have unbounded maximum velocity at $T$. Then
\begin{equation}\label{eq:closure}
 \mathcal R_{a,T}\subseteq\overline{\mathcal S_{a,T}}^{\,E_T}.
\end{equation}
For each reference pair $(v,g)$ and every fixed $s<1/2$, the approximating pairs can be chosen with
\[
 (u_\eps,g_\eps)\longrightarrow(v,g)
 \quad\hbox{in }E_T\times L^1(0,\infty;H^s).
\]
\end{corollary}
\begin{proof}
A reference smooth through $T$ is smooth through $T+\delta$ for some
$\delta>0$, by local continuation; equivalently one may take this as the meaning of smoothness through the endpoint. Apply Theorem~\ref{thm:insertion}. Its energy estimate gives convergence in $E_T$, and its force estimates give simultaneous convergence. The reference has finite $E_T$ norm and the packet and correction do also, so $u_\eps$ belongs to the specified ambient energy space.

To verify the claimed interpretation of the norm, note that
\[
 \norm{z}_{L^2(0,T;L^2)}
 \le T^{1/2}\norm{z}_{L^\infty(0,T;L^2)}.
\]
Thus $E_T$ controls both terms of $L^2_tH^1_x$, while the usual sum norm controls $E_T$ directly. This proves the equivalence used in \eqref{eq:closure}. No endpoint value or classical continuation of $u_\eps$ is implied by convergence in this norm.
\end{proof}

\begin{proposition}[Projection of extended singular data]\label{prop:projection}
Define
\[
 \mathfrak B_{\nu,T}
 =\{(a,f)\in\XX\times\FF:T_{\max}^{\nu}(a,f)\le T\}.
\]
If $s<1/2$, this set is dense in the product of any topology on $\XX$ and the relative $L^1_tH^s_x$ topology on $\FF$. Its projection onto $\XX$ is all of $\XX$.
The family obtained while requiring $a=0$ projects instead to the singleton $\{0\}$.
\end{proposition}
\begin{proof}
Take a point $(a,g)$ and a basic product neighborhood $U\times V$ containing it. The initial datum $a$ already belongs to $U$. Proposition~\ref{prop:density} supplies a force $f\in V$ with $T_{\max}^{\nu}(a,f)\le T$. Hence
$(a,f)\in\mathfrak B_{\nu,T}\cap(U\times V)$, proving density.
Applying the same proposition with any fixed $g$ and any positive radius gives at least one singular force over each $a$, which proves the projection statement. The final assertion follows directly from fixing the initial velocity to zero.
\end{proof}

The projection statement has the quantifier order
$\forall a\,\exists f$, not $\exists f\,\forall a$.
For a fixed prescribed force $f_*$, the set
\[
 \{a\in\XX:T_{\max}^{\nu}(a,f_*)\le T\}
\]
is a different object and is not classified by these force-changing constructions.
In a nontrivial normed initial-velocity space, $\{0\}$ is not dense: if
$a\ne0$, the ball of radius $\norm{a}/2$ centered at $a$ excludes zero.

There is also a distinction between the two terminal-time assertions.
Theorem~\ref{thm:insertion} gives singularity exactly at $T$ around every reference regular through $T$. Proposition~\ref{prop:density} uses breakdown by $T$, because a general reference may already break down earlier and is then left unchanged. It does not prove density of exactly-time-$T$ singular forces around every earlier-breaking reference.

\section{Critical forcing and a regular open neighborhood}\label{sec:critical}
The positive result does not decide what happens at $s=1/2$. We now prove a regularity estimate that is independent of the packet construction. The critical Sobolev exponent is consistent with the classical Fujita--Kato framework \cite{fujita}; the forcing, mean mode and continuation steps needed here are included explicitly.

\begin{proposition}[Global regularity for small critical force]\label{prop:critical}
There is $c>0$, depending only on the unit torus and the stated norm conventions, such that
\begin{equation}\label{eq:smallcritical}
 g\in\FF,\qquad
 \rho:=\norm{g}_{L^1(0,\infty;H^{1/2})}<c\nu
\end{equation}
imply $T_{\max}^{\nu}(0,g)=\infty$.
\end{proposition}
\begin{proof}
Let $u$ be the maximal classical solution from Proposition~\ref{prop:local}.
All estimates are first made on a compact subinterval of its lifespan.

\emph{Removal of the spatial mean.}
Define
\[
 \overline g(t)=\int_{\TT}g(x,t)\dd x,\qquad
 m(t)=\int_{\TT}u(x,t)\dd x,\qquad
 v=u-m,\qquad h=g-\overline g.
\]
Integrating the periodic equation eliminates the convective, Laplacian and pressure terms, giving
$m'=\overline g$ and $m(0)=0$. Thus
\begin{equation}\label{eq:meanbound}
 |m(t)|\le\int_0^t|\overline g(s)|\dd s\le\rho.
\end{equation}
The mean-zero field $v$ solves
\begin{equation}\label{eq:meanfree}
 \partial_tv+(v\cdot\nabla)v+(m\cdot\nabla)v-\nu\Delta v+\nabla p=h.
\end{equation}
For any real Sobolev order, the spatially constant transport $m(t)\cdot\nabla$ commutes with the corresponding Fourier multiplier and is skew-adjoint in $L^2$. It therefore contributes zero to both energy identities used below. This step is necessary because the force is not assumed to have zero spatial mean.

\emph{Critical energy estimate.}
Let $\Lambda=(-\Delta)^{1/2}$ on mean-zero fields, and set
\[
 y(t)=\norm{\Lambda^{1/2}v(t)}_2,\qquad
 z(t)=\norm{\Lambda^{3/2}v(t)}_2,\qquad
 b(t)=\norm{h(t)}_{\dot H^{1/2}}.
\]
Testing \eqref{eq:meanfree} against $\Lambda v$ gives
\[
 \frac12(y^2)'+\nu z^2
 =-\ip{(v\cdot\nabla)v}{\Lambda v}+\ip{h}{\Lambda v}.
\]
The pressure term is zero because $\nabla\cdot\Lambda v=0$.
By Lemma~\ref{lem:calculus} and H\"older's inequality,
\[
 |\ip{(v\cdot\nabla)v}{\Lambda v}|
 \le\norm{v}_3\norm{\nabla v}_3\norm{\Lambda v}_3
 \le C_0yz^2.
\]
The force term is bounded by
\[
 |\ip{h}{\Lambda v}|
 =|\ip{\Lambda^{1/2}h}{\Lambda^{1/2}v}|\le b(t)y(t).
\]
Consequently,
\begin{equation}\label{eq:criticalenergy}
 \frac12(y^2)'+(\nu-C_0y)z^2\le by.
\end{equation}
Removing the zero Fourier mode cannot increase the relevant norm, so
\begin{equation}\label{eq:bintegral}
 \int_0^\infty b(t)\dd t\le\rho.
\end{equation}

\emph{Closing the bootstrap, including the zero norm.}
On any interval where $y\le\nu/(2C_0)$, \eqref{eq:criticalenergy} implies
$\tfrac12(y^2)'\le by$. For $\zeta>0$,
\[
 \frac{\dd}{\dd t}(y^2+\zeta^2)^{1/2}
 \le b\,\frac{y}{(y^2+\zeta^2)^{1/2}}\le b.
\]
Since $y(0)=0$, integration and $\zeta\downarrow0$ yield
\begin{equation}\label{eq:ybound}
 y(t)\le\int_0^t b(s)\dd s\le\rho
\end{equation}
throughout such an interval. Choose $c<1/(4C_0)$.
If a first time with $y=\nu/(2C_0)$ existed, continuity would permit the preceding calculation up to that time, where it would give the strictly smaller bound $y\le\rho<\nu/(4C_0)$. This contradiction closes the bootstrap on the whole classical lifespan.

In particular the critical dissipation also obeys
\[
 \frac12y(t)^2+\frac{\nu}{2}\int_0^t z(s)^2\dd s
 \le\int_0^t b(s)y(s)\dd s\le\rho^2.
\]
We next obtain a stronger integrability bound that directly implies classical continuation.

\emph{The $H^1$ estimate.}
Take the inner product of \eqref{eq:meanfree} with $-\Delta v$.
The pressure and mean-transport terms vanish. The nonlinear term satisfies
\begin{align*}
 |\ip{(v\cdot\nabla)v}{\Delta v}|
 &\le\norm{v}_3\norm{\nabla v}_6\norm{\Delta v}_2\\
 &\le C_1y\,\norm{\Delta v}_2^2.
\end{align*}
Here $\norm{\nabla v}_6\le C\norm{\Delta v}_2$ follows from
Lemma~\ref{lem:critical-embeddings}, applied to each mean-zero derivative:
$\norm{\nabla(\partial_jv)}_2\le\norm{\Delta v}_2$ by its Fourier series. Choose also $c<1/(4C_1)$.
By \eqref{eq:ybound}, the convection term is at most
$(\nu/4)\norm{\Delta v}_2^2$. Young's inequality bounds the force term by
$(\nu/4)\norm{\Delta v}_2^2+C\nu^{-1}\norm{h}_2^2$.
After multiplying the resulting inequality by two and weakening its positive coefficient if necessary,
\begin{equation}\label{eq:H1energy}
 (\norm{\nabla v}_2^2)'
 +\nu\norm{\Delta v}_2^2\le C\nu^{-1}\norm{h}_2^2.
\end{equation}
Since $v(0)=0$, integration gives a uniform bound on
$\norm{\nabla v(t)}_2^2$ and on $\int_0^t\norm{\Delta v(s)}_2^2\dd s$.
The integral of $\norm{h}_2^2$ over all positive times is finite because each fixed force is smooth and time compact. It need not be small.

\emph{Continuation.}
Suppose the maximal lifespan were a finite number $S$.
For mean-zero $v$, Fourier series give
$\norm{v}_{H^2}\le C\norm{\Delta v}_2$.
Together with \eqref{eq:H1energy} and \eqref{eq:meanbound}, the
orthogonality of the constant and mean-zero Fourier modes gives
\begin{align*}
 \int_0^S\norm{u(t)}_{H^2}^2\dd t
 &=\int_0^S\bigl(|m(t)|^2+\norm{v(t)}_{H^2}^2\bigr)\dd t\\
 &\le S\rho^2+C\nu^{-2}
             \int_0^\infty\norm{h(t)}_2^2\dd t<\infty.
\end{align*}
Proposition~\ref{prop:local} then extends $u$ beyond $S$, a contradiction.
Thus the maximal lifespan is infinite. The choice of $c$ uses only the constants in the periodic Sobolev estimates, not the amplitude, derivatives or support size of a particular force.
\end{proof}

\begin{corollary}[Non-density at and above the critical exponent]\label{cor:nondensity}
For every $s\ge1/2$ and $T>0$, $\BB^0_{\nu,T}$ is not dense in
$\FF$ with the relative $L^1_tH^s_x$ topology.
\end{corollary}
\begin{proof}
At $s=1/2$, the set
\[
 \{g\in\FF:\norm{g}_{L^1_tH^{1/2}_x}<c\nu\}
\]
is a nonempty relative open ball, containing zero, and is disjoint from
$\BB^0_{\nu,T}$ by Proposition~\ref{prop:critical}.
For $s\ge1/2$, the Fourier weights imply
$\norm{g(t)}_{H^{1/2}}\le\norm{g(t)}_{H^s}$.
The ball $\norm{g}_{L^1_tH^s_x}<c\nu$ is therefore also disjoint from the singular-force set. A dense subset cannot miss a nonempty open set.
\end{proof}

\begin{proof}[Proof of Theorem~\ref{thm:main}]
Proposition~\ref{prop:density} gives density for every fixed smooth initial velocity when $s<1/2$. Corollary~\ref{cor:nondensity} gives non-density on the zero-data slice when $s\ge1/2$. These are the two assertions of the theorem.
\end{proof}
The non-density statement is proved for the zero-data slice.
It does not classify every nonzero-data slice at or above the critical exponent. In contrast, the positive insertion preserves any fixed smooth initial velocity.

\section{Additional singular families}\label{sec:variants}
\begin{proposition}[Infinite-dimensional variations away from the endpoints]\label{prop:affine}
Fix the packet $(U,P,F)$ of Theorem~\ref{thm:packet}, and fix a spacetime cylinder $Q=B_0\times(\tau_0,\tau_1)$ with
$0<\tau_0<\tau_1<1$ and $B_0\Subset\R^3$ an open ball.
For any $b\in C_c^\infty(Q;\R^3)$ with $\nabla\cdot b=0$, define
\begin{align}
 \widetilde U&=U+b,\qquad \widetilde P=P,\nonumber\\
 \widetilde F&=F+\partial_tb-\nu\Delta b
       +(U\cdot\nabla)b+(b\cdot\nabla)U+(b\cdot\nabla)b.\label{eq:affine}
\end{align}
These fields give an infinite-dimensional affine family of distinct velocities with zero initial data, finite energy and dissipation, and the same late singular behavior as $U$. For each fixed nonzero $b$, the family obtained by replacing $b$ with $\lambda b$ is non-isolated as $\lambda\to0$ in every fixed smooth seminorm of the velocity and force differences on their compact support.
\end{proposition}
\begin{proof}
Expanding the momentum equation for $U+b$ gives exactly the force in
\eqref{eq:affine}; its divergence remains zero. Every correction term is supported in $\supp b$, a compact subset of $Q$. The original packet is smooth with bounded derivatives of every fixed order on a neighborhood of that set. Thus the force correction extends smoothly by zero to the whole positive time axis, including time one, despite the singularity of $U$ there.

Since $b=0$ near time zero and near time one,
$\widetilde U(0)=0$ and $\widetilde U=U$ on a fixed late interval.
The same unbounded-speed limit holds. The energy and dissipation norms of $b$ are finite by compact smooth support, so the triangle inequality gives the corresponding bounds for $\widetilde U$.

To verify infinite dimensionality explicitly, choose countably many disjoint small balls inside $B_0$. In each ball choose a smooth compactly supported vector potential with nonzero curl, and multiply that curl by a fixed nonzero time bump in $(\tau_0,\tau_1)$. These divergence-free fields are linearly independent because their spatial supports are disjoint. Their finite linear combinations belong to the allowed class. The map $b\mapsto U+b$ is injective and affine, so its image is infinite dimensional.

For the non-isolation statement, the velocity difference is $\lambda b$, while the force difference has the form
\[
 \lambda L_U b+\lambda^2(b\cdot\nabla)b,\qquad
 L_Ub=\partial_tb-\nu\Delta b+(U\cdot\nabla)b+(b\cdot\nabla)U.
\]
On the fixed compact support, all coefficients and their derivatives are bounded. Hence for each integer $m\ge0$ the $C^m$ norm of the force difference is at most $C_m|\lambda|+C_m'\lambda^2$, and the velocity difference is at most $C_m''|\lambda|$. Both tend to zero. The forces vary in this family; no open singularity basin for one fixed force follows.
\end{proof}

\begin{proposition}[Finitely many prescribed singular regions]\label{prop:multiple}
Fix finitely many disjoint interior balls $B_1,\ldots,B_N$ in the torus or in a bounded three-dimensional domain. At fixed $\nu>0$ and $T>0$, there is a smooth forced solution starting from rest, smooth on $[0,T)$, with finite energy and dissipation, such that
\[
 \limsup_{t\uparrow T}\norm{u(t)}_{L^\infty(B_j)}=\infty
 \qquad(1\le j\le N).
\]
In a bounded domain its boundary condition is homogeneous no-slip.
\end{proposition}
\begin{proof}
Choose $x_j\in B_j$ and scales $\eps_j>0$ sufficiently small that the complete scaled compact set $x_j+\eps_jK_*$ lies strictly inside $B_j$ and $\eps_j^2<T$. Apply \eqref{eq:scaling} to each packet with start time $T-\eps_j^2$, and denote the resulting fields by $U_j,P_j,F_j$.
Set
\[
 u=\sum_{j=1}^N U_j,\qquad p=\sum_{j=1}^N P_j,\qquad f=\sum_{j=1}^N F_j.
\]
For $i\ne j$, the supports of $U_i$ and $U_j$ are separated, so
$(U_i\cdot\nabla)U_j=0$. The sum therefore satisfies the momentum equation exactly. All force supports are compact in positive time and all forces are smooth. Initial vanishing gives $u(0)=0$.

On $B_j$ the velocity equals $U_j$, proving the stated separate limsup for each ball. Disjoint supports also give
\[
 \sup_{t<T}\norm{u(t)}_2^2\le M^2\sum_{j=1}^N\eps_j,\qquad
 \int_0^T\norm{\nabla u(t)}_2^2\dd t=D^2\sum_{j=1}^N\eps_j.
\]
These sums are finite. In a bounded domain, all fields vanish in a boundary collar; hence $u=0$ on the boundary. On the torus a spatially constant pressure adjustment gives the chosen gauge.
\end{proof}
The proposition concerns a finite number of packets with a common terminal time. It asserts neither an infinite superposition nor a classical continuation after the first singular endpoint. The times at which individual packet maxima are large need not coincide; the displayed conclusions are separate limsup statements as $t\uparrow T$.

\begin{proposition}[Conservative forcing from rest]\label{prop:conservative}
On the torus, let $f=-\nabla\phi$ for a globally defined periodic scalar
$\phi$. On a bounded domain, let $f=-\nabla\phi$ and impose homogeneous no-slip velocity. In either case, any smooth solution starting from rest is identically zero on its classical lifespan.
\end{proposition}
\begin{proof}
Incompressibility and periodicity or the zero normal boundary trace give
\[
 \int f\cdot u
 =-\int\nabla\phi\cdot u=0.
\]
The classical energy identity, integrated from zero, becomes
\[
 \frac12\norm{u(t)}_2^2
 +\nu\int_0^t\norm{\nabla u(s)}_2^2\dd s=0.
\]
Both terms are nonnegative, so $u(t)=0$ at every presingular time.
The periodic-potential condition excludes affine potentials, whose gradients need not be removable by the allowed periodic pressure gauge.
\end{proof}

Every nontrivial from-rest construction in this article therefore requires a nonconservative force. Smoothness of the zero initial velocity does not imply that the resulting force-velocity pair is regular through the terminal time.

\section{Discussion}\label{sec:discussion}
The main theorem gives a complete dense/not-dense classification in a particular scale of force topologies on the zero-initial-velocity slice. Below $s=1/2$, the force norm permits enough concentration for a singular packet to be inserted near any regular reference. At and above the critical exponent, an independent estimate provides an open set of globally regular smooth forces. The critical conclusion is an analytic obstruction, not an inference drawn solely from scaling.

The trajectory result is equally specific. Singular solutions can approach a regular reference strongly in both the kinetic-energy and integrated-dissipation terms of $E_T$. At each presingular time they are smooth and belong to every finite-order spatial Sobolev space. These facts do not bound their maximum velocity uniformly as the terminal time is approached.

The extended-data projection reaches every smooth initial velocity because the force is free to depend on that velocity. This does not determine the density, dimension or stability of singular initial states under a single prescribed force. The literal from-rest family has only the singleton initial datum. Maintaining these quantifiers is essential to any interpretation of the distribution result.

No probability law on forces or perturbations has been specified. Density does not establish typicality, an open basin of singularity, or robustness under arbitrary errors. Likewise, density in an infinite-dimensional input space does not imply density within a fixed spectral truncation or actuator family. The approximating forces have shorter scales and increasing peaks. A numerical conclusion must therefore identify the retained observations and the error norm.

The results do not select a continuation after the singular time. No CFD experiment is claimed.

\section*{Declaration of generative AI use}
GPT Astra (OpenAI) was used to assist with literature and Lean-source inspection, development and checking of mathematical arguments, adversarial review, drafting, and LaTeX preparation. The use of generative AI does not constitute formal verification. Responsibility for the final content remains with the authors.

\appendix

\section{A direct proof of the critical embeddings}\label{sec:critical-appendix}
This appendix proves the two critical embeddings used in the nonlinear
estimates.  Only elementary integration, Fourier inversion and Plancherel,
H\"older's inequality, and approximation in $L^2$ are used.  In particular,
no Hardy--Littlewood--Sobolev or Sobolev embedding theorem is assumed.
Throughout this appendix $\Lambda=(-\Delta)^{1/2}$, with Fourier symbol
$|\xi|$ on $\R^3$ and $2\pi|k|$ on $\TT=\R^3/\mathbb Z^3$.

\begin{lemma}[Critical embeddings on the two domains]\label{lem:critical-embeddings}
For $a\in\{1/2,1\}$ and $p_a=6/(3-2a)$, there are finite constants,
depending only on $a$ and on the fixed unit torus when relevant, such that
\begin{equation}\label{eq:critical-embedding-pair}
 \norm{v}_{L^{p_a}(\R^3)}\le C_a\norm{\Lambda^a v}_{L^2(\R^3)},
 \qquad
 \norm{v}_{L^{p_a}(\TT)}\le C_a'\norm{\Lambda^a v}_{L^2(\TT)}.
\end{equation}
The first inequality holds initially for Schwartz functions and extends
to the homogeneous completion specified below.  The second holds for
mean-zero periodic distributions with finite displayed right-hand side;
the mean-zero condition is essential.  Componentwise application gives
the same statements for vector and tensor fields, with constants allowed
to depend on their fixed number of components.

Consequently, on either domain and whenever the indicated homogeneous
norms are finite (with mean-zero representatives on the torus),
\begin{align}
 \norm{v}_3&\le C\norm{v}_{\dot H^{1/2}},\nonumber\\
 \norm{\nabla v}_3+\norm{\Lambda v}_3
   &\le C\norm{v}_{\dot H^{3/2}},\label{eq:critical-derived}\\
 \norm{\nabla v}_6&\le C\norm{\Delta v}_2.\nonumber
\end{align}
Here the derivative statements on $\R^3$ refer to distributions whose
derivatives have the homogeneous realizations described in the proof;
in particular they apply to all the smooth $H^\infty$ fields used in
the article.
\end{lemma}

\begin{proof}
\emph{1. A maximal-function estimate, including its proof.}
For a locally integrable $g$ on $\R^3$ let
\[
 Mg(x)=\sup_{r>0}\frac{1}{|B(x,r)|}
                    \int_{B(x,r)}|g(y)|\dd y.
\]
For $g\in L^1$ and $\lambda>0$, the set $E_\lambda=\{Mg>\lambda\}$
is open: for each fixed radius the ball integral is continuous in its
center, by translation continuity in $L^1$ on bounded sets.  Every compact
$K\subset E_\lambda$ is covered by finitely many balls whose averages
exceed $\lambda$.  Select the largest remaining ball and discard all balls
meeting it, then repeat.  The selected balls $B_j$ are disjoint, and each
discarded ball lies in the concentric triple of the ball that discarded it.
Therefore
\[
 |K|\le\sum_j|3B_j|=27\sum_j|B_j|
       \le\frac{27}{\lambda}\norm{g}_1.
\]
Exhaustion by compact subsets proves
$|E_\lambda|\le27\lambda^{-1}\norm{g}_1$.

For $g\in L^2$, split $g=g\mathbf1_{\{|g|\le\lambda/2\}}
+g\mathbf1_{\{|g|>\lambda/2\}}$.  The first part has maximal function at
most $\lambda/2$, and the second is in $L^1$.  The preceding estimate gives
\[
 |\{Mg>\lambda\}|\le\frac{54}{\lambda}
             \int_{\{|g|>\lambda/2\}}|g(x)|\dd x.
\]
The layer-cake identity and Tonelli's theorem now yield
\begin{equation}\label{eq:critical-maximal-ltwo}
 \norm{Mg}_2^2
 =2\int_0^\infty\lambda|\{Mg>\lambda\}|\dd\lambda
 \le108\int_{\R^3}|g(x)|
             \int_0^{2|g(x)|}\dd\lambda\dd x
 =216\norm{g}_2^2.
\end{equation}
This also proves that $Mg$ is finite almost everywhere; no prior
integrability of $Mg$ was needed.

\emph{2. The potential kernel from the Fourier multiplier.}
For $0<a<3$, define $I_a$ by the multiplier $|\xi|^{-a}$ on Schwartz
functions.  The Gamma integral gives, for $\xi\ne0$,
\[
 |\xi|^{-a}=\frac{1}{\Gamma(a/2)}
       \int_0^\infty t^{a/2-1}e^{-t|\xi|^2}\dd t,
 \qquad
 \Gamma(b)=\int_0^\infty s^{b-1}e^{-s}\dd s.
\]
Direct Gaussian Fourier integration shows that the heat multiplier
$e^{-t|\xi|^2}$ acts by convolution with
$h_t(x)=(4\pi t)^{-3/2}e^{-|x|^2/(4t)}$ in the unitary Fourier convention
of this article.  Substituting $s=|x|^2/(4t)$ gives
\begin{align*}
 K_a(x)&=\frac{1}{\Gamma(a/2)}
           \int_0^\infty t^{a/2-1}h_t(x)\dd t\\
 &=\frac{\Gamma((3-a)/2)}{2^a\pi^{3/2}\Gamma(a/2)}|x|^{a-3}
   =:c_a|x|^{a-3},\qquad x\ne0.
\end{align*}
Thus $I_ag=K_a*g$.  To justify the identity without a formal interchange,
for Schwartz $g$ the small-$t$ integrand is bounded by
$t^{a/2-1}\norm{g}_\infty$, and the large-$t$ integrand is bounded by
$C t^{(a-3)/2-1}\norm{g}_1$.  Both are integrable in their respective
ranges.  The same bounds after pairing with a Schwartz test function
justify the heat representation in distributions.  On the Fourier side,
$|\xi|^{-a}$ is locally integrable for $a<3$ and the other factors decay
rapidly, so Fubini gives exactly the displayed multiplier.  On the spatial
side the locally integrable kernel and the decay of $g$ justify the
convolution and its heat-integral representation.

\emph{3. The potential estimate.}
Take $0<a<3/2$, $g\in L^2$ and $R>0$.  Splitting the potential at radius
$R$, and dividing the inner ball into the annuli with outer radii
$2^{-j}R$, gives
\begin{align*}
 \int_{|y|<R}|y|^{a-3}|g(x-y)|\dd y
 &\le C_a\sum_{j=0}^\infty(2^{-j}R)^a Mg(x)
 \le C_a R^a Mg(x),\\
 \int_{|y|\ge R}|y|^{a-3}|g(x-y)|\dd y
 &\le\norm{g}_2
       \left(4\pi\int_R^\infty r^{2a-4}\dd r\right)^{1/2}
 \le C_a\norm{g}_2 R^{a-3/2}.
\end{align*}
The first inequality uses the average over each outer ball and the maximum
of the kernel on that annulus.  The second integral converges precisely
because $a<3/2$.  At points where $0<Mg(x)<\infty$ choose
$R=(\norm{g}_2/Mg(x))^{2/3}$.  The zero function is harmless; if $Mg(x)=0$
at some point, the integrals over every ball centered there show that
$g=0$ almost everywhere.  We have proved almost everywhere
\[
 |I_ag(x)|\le C_a\norm{g}_2^{2a/3}Mg(x)^{1-2a/3}.
\]
This also proves absolute convergence of the kernel integral almost
everywhere for $g\in L^2$.  With $p_a=6/(3-2a)$, the identity
$p_a(1-2a/3)=2$ and \eqref{eq:critical-maximal-ltwo} imply
\begin{equation}\label{eq:critical-potential}
 \norm{I_ag}_{p_a}
 \le C_a\norm{g}_2^{2a/3}\norm{Mg}_2^{1-2a/3}
 \le C_a\norm{g}_2.
\end{equation}

\emph{4. The Euclidean function class and passage to the limit.}
First let $v$ have Fourier transform in
$C_c^\infty(\R^3\setminus\{0\})$.  Then $g=\Lambda^a v$ is Schwartz,
the multiplier identity in step 2 gives $v=I_ag$, and
\eqref{eq:critical-potential} proves the Euclidean inequality.
For clarity, the homogeneous space here is the completion of this class
in $\norm{\Lambda^a v}_2$, for $0<a<3/2$.  It has a unique realization as
distributions, and the estimate gives a unique $L^{p_a}$ representative.
Indeed, if $G\in L^2$ is the limit of $|\xi|^a\widehat v_n$, its
distributional representative has Fourier transform $|\xi|^{-a}G$:
for every Schwartz test function $\phi$,
\[
 \int_{\R^3}|\xi|^{-a}|G(\xi)\widehat\phi(\xi)|\dd\xi
 \le\norm{G}_2
      \left(\int_{\R^3}|\xi|^{-2a}|\widehat\phi(\xi)|^2\dd\xi\right)^{1/2}
 <\infty.
\]
The last integral is finite at zero because $2a<3$, and finite at infinity
by rapid decay.  This estimate also proves distributional convergence.
Meanwhile the embedding makes $v_n$ Cauchy in $L^{p_a}$, and convergence
in that space implies distributional convergence by H\"older's inequality;
the two limits coincide.  Every $G\in L^2$ is approximable by smooth
functions supported in compact annuli: first truncate in frequency, then
mollify within a slightly larger annulus.  Hence this realization covers
every such $G$, without an ambiguity by additive polynomials.  In
particular any Schwartz $v$, and any $H^\infty$ field with the stated
homogeneous norm, is covered by approximating
$G=|\xi|^a\widehat v$ in this way.

\emph{5. Transfer to the unit torus.}
Choose $\eta\in C_c^\infty(\R^3)$ equal to one on the cube
$[-1/2,1/2]^3$.  For a trigonometric polynomial $v$, extend it periodically
and set $V=\eta v$.  Write $\langle\xi\rangle=(1+|\xi|^2)^{1/2}$.
Fourier transformation gives
\[
 \widehat V(\xi)=\sum_{k\in\mathbb Z^3}\widehat v(k)
                                  \widehat\eta(\xi-2\pi k).
\]
Integration by parts in the compactly supported smooth function $\eta$
gives arbitrarily fast decay of $\widehat\eta$.  Together with
$\langle\xi\rangle^a\le C_a\langle2\pi k\rangle^a
\langle\xi-2\pi k\rangle^a$, this yields
\[
 \langle\xi\rangle^a|\widehat V(\xi)|
 \le C_{a,\eta}\sum_k A_k Q(\xi-2\pi k),
 \quad
 A_k=\langle2\pi k\rangle^a|\widehat v(k)|,
 \quad Q(z)=(1+|z|^2)^{-2}.
\]
Here $Q\in L^1(\R^3)$, and
$\sup_\xi\sum_k Q(\xi-2\pi k)<\infty$: reduce $\xi$ to a fixed
lattice cell and group the lattice points into shells, whose contributions
are bounded by $C\sum_{n\ge0}(n+1)^2(1+n^2)^{-2}<\infty$.
Weighted Cauchy--Schwarz, followed by integration, therefore proves
\begin{align*}
 \norm{V}_{H^a(\R^3)}^2
 &\le C_{a,\eta}\int_{\R^3}
       \left(\sum_k A_k^2Q(\xi-2\pi k)\right)
       \left(\sum_k Q(\xi-2\pi k)\right)\dd\xi\\
 &\le C_{a,\eta}\sum_k A_k^2
  =C_{a,\eta}\norm{v}_{H^a(\TT)}^2.
\end{align*}
Apply the Euclidean estimate to $V$ and restrict to the unit cell.  If
$\widehat v(0)=0$, the elementary spectral-gap estimate
\[
 (1+4\pi^2|k|^2)^a
 \le(1+(2\pi)^{-2})^a(2\pi|k|)^{2a},\qquad k\ne0,
\]
then gives
\[
 \norm{v}_{L^{p_a}(\TT)}
 \le\norm{V}_{L^{p_a}(\R^3)}
 \le C_a\norm{V}_{\dot H^a(\R^3)}
 \le C_a'\norm{v}_{\dot H^a(\TT)}.
\]
Fourier truncations converge in $\dot H^a$ for every mean-zero periodic
distribution with finite norm, hence also in $L^{p_a}$ by this estimate.
Their distributional limits identify the two representatives.  This proves
the periodic statement, including its extension beyond polynomials.

Finally apply the $a=1/2$ estimate to each $\partial_jv$ and to $\Lambda v$,
and the $a=1$ estimate to each $\partial_jv$.  Plancherel and
$\sum_j|\xi_j|^2=|\xi|^2$ (or its Fourier-series version) give
\eqref{eq:critical-derived}.  For smooth fields the resulting nonlinear
bound is, explicitly,
\[
 |\ip{(v\cdot\nabla)v}{\Lambda v}|
 \le\norm{v}_3\norm{\nabla v}_3\norm{\Lambda v}_3
 \le C\norm{v}_{\dot H^{1/2}}\norm{v}_{\dot H^{3/2}}^2.
\]
Every constant above is independent of the field and of its frequency
support.
\end{proof}

\begin{thebibliography}{9}\small\raggedright
\bibitem{openai}
OpenAI. \emph{Finite Time Blowup for Navier--Stokes}. Manuscript, 166 pp., accessed 9 September 2026.
\href{https://cdn.openai.com/pdf/32d9f210-8b73-45e0-91bc-82a30aef8a9a/navier-stokes.pdf}{Official manuscript}.
\bibitem{lean}
OpenAI. \emph{NavierStokesAndEuler}. Lean source repository, revision
\texttt{8937a8f4cbc7abaab5e9e97d1cc7f5d2319d9538}, accessed 9 September 2026.
\href{https://github.com/openai/NavierStokesAndEuler/tree/8937a8f4cbc7abaab5e9e97d1cc7f5d2319d9538}{Source revision}.
Compact witness: \href{https://github.com/openai/NavierStokesAndEuler/blob/8937a8f4cbc7abaab5e9e97d1cc7f5d2319d9538/NavierStokes/R3ActualCandidate.lean#L18-L21}{\texttt{selected\_compact\_candidate}};
its support and zero-data conditions: \href{https://github.com/openai/NavierStokesAndEuler/blob/8937a8f4cbc7abaab5e9e97d1cc7f5d2319d9538/NavierStokes/R3CompactCandidate.lean#L24-L37}{\texttt{R3CompactCandidate.Properties}}.
\bibitem{hzz2024}
M. Hofmanov\'a, R. Zhu and X. Zhu.
\emph{Non-uniqueness of Leray--Hopf solutions for stochastic forced Navier--Stokes equations}.
arXiv:2309.03668v2, 10 September 2024.
\href{https://arxiv.org/abs/2309.03668}{arXiv:2309.03668}.
\bibitem{fujita}
H. Fujita and T. Kato. On the Navier--Stokes initial value problem. I.
\emph{Archive for Rational Mechanics and Analysis} \textbf{16} (1964), 269--315.
\href{https://doi.org/10.1007/BF00276188}{doi:10.1007/BF00276188}.
\bibitem{tao2013}
T. Tao. Localisation and compactness properties of the Navier--Stokes global regularity problem.
\emph{Analysis \& PDE} \textbf{6}(1) (2013), 25--107.
\href{https://doi.org/10.2140/apde.2013.6.25}{doi:10.2140/apde.2013.6.25}.
\bibitem{tao2018}
T. Tao. \emph{254A, Notes 1: Local well-posedness of the Navier--Stokes equations}.
Author's course notes, 16 September 2018.
\href{https://terrytao.wordpress.com/2018/09/16/254a-notes-1-local-well-posedness-of-the-navier-stokes-equations/}{Original text}.
\end{thebibliography}
\end{document}
